%!TeX encoding = UTF-8 Unicode
%!TeX spellcheck = en-UK
% !TeX root = Perrinet26mnesis.tex
% -*- coding: UTF-8; -*-
%%%-----------------------------------------------------------------
\begin{figure}[!t] 
  \centering
  \includegraphics[width=0.33\linewidth]{scan_num_delay.pdf}\hfill
  \includegraphics[width=0.33\linewidth]{scan_N_time.pdf}\hfill
  \includegraphics[width=0.33\linewidth]{scan_N_neuron.pdf}
  % TODO \includegraphics[width=0.9\linewidth]{scan_p_A.pdf}
  \caption{%
    Effect of some key parameters on working-memory performance (loss $\mathcal{L} = 1 - F_1$, mean $\pm$ s.e.\ over $N_\mathrm{cv} = 10$ cross-validation folds). %
    \emph{Left:} Increasing the maximum delay $D$ consistently reduces the loss, confirming that deeper delay windows provide a larger, more orthogonal context and that longer maximum delays directly increase the capacity of the recurrent HD-SNN. % (Section~\ref{sec:capacity}). %
    \emph{Middle:} 
    Loss grows slowly and monotonically with pattern duration $T$, consistent with the compounding-error and gradient-vanishing arguments: longer sequences require the network to maintain coherent predictions further from the trigger window. %
    \emph{Right:} 
    Performance is optimal for a number of neurons superior to $\approx 500$ the loss deteriorating rapidly for lower numbers of neurons.
    % Performance is optimal at low background firing rates ($p_A \lesssim 10^{-3}$) and degrades sharply at higher rates. 
  }\label{fig:scan}
\end{figure}
%%%-----------------------------------------------------------------
